\documentclass[a4paper]{article}
\usepackage[utf8]{inputenc}
\usepackage[T2A]{fontenc}
\usepackage[english,russian]{babel}
\usepackage[left=25mm, top=20mm, right=25mm, bottom=30mm, nohead, nofoot]{geometry}
\usepackage{amsmath,amsfonts,amssymb}
\usepackage{fancybox,fancyhdr}
\usepackage{enumitem}
\usepackage{xcolor}
\usepackage{hyperref}

\hypersetup{colorlinks=true, allcolors=[RGB]{010 090 200}}

\pagestyle{fancy}
\fancyhf{}
\fancyhead[R]{Информатика ЕГЭ — Задание 16}
\fancyfoot[R]{\thepage}
\headsep=10mm

\definecolor{easycolor}{RGB}{0,150,0}
\definecolor{medcolor}{RGB}{200,100,0}
\definecolor{hardcolor}{RGB}{180,0,0}

\newcommand{\easytag}{\textcolor{easycolor}{\textbf{[easy]}}}
\newcommand{\medtag}{\textcolor{medcolor}{\textbf{[medium]}}}
\newcommand{\hardtag}{\textcolor{hardcolor}{\textbf{[hard]}}}
\newcommand{\answer}[1]{\par\smallskip\textit{\textbf{Ответ:} #1}}
\newcommand{\solution}[1]{\par\textit{\textbf{Решение:} #1}}

% ============================================================
% КАК ДОБАВИТЬ НОВУЮ ЗАДАЧУ:
% \item % ID: cs_ege_16_???_00N | \???tag | src: №XXXXX
% Условие...
% \answer{...}
% \solution{...}
% ============================================================

\begin{document}

\Large
\section*{Задание 16 — Рекурсивные алгоритмы}

\normalsize
\textit{Тема:} рекурсивное определение функций, вычисление значений по соотношениям.\\
\textit{Источник:} \href{https://inf-ege.sdamgia.ru/}{СдамГИА — Информатика ЕГЭ}

\bigskip

% ========================
%  ЛЁГКИЕ ЗАДАЧИ (easy)
% ========================

\subsection*{\textcolor{easycolor}{Лёгкий уровень}}

\begin{enumerate}[label=\textbf{Задача \arabic*.}, leftmargin=*, itemsep=10mm]

  \item % ID: cs_ege_16_easy_001 | \easytag | src: №7308
  \easytag\ Алгоритм вычисления значения функции $F(n)$, где $n$ — натуральное число,
  задан следующими соотношениями:
  \[
    F(1) = 1;\quad F(n) = F(n-1) + 2n - 1,\quad n > 1.
  \]
  Чему равно значение функции $F(10)$?
  \textit{В ответе запишите только натуральное число.}
  \answer{100}
  \solution{$F(n) = n^2$. Проверка: $F(1)=1,\ F(2)=1+3=4,\ F(3)=4+5=9,\ldots,\ F(10)=100$.}

  \bigskip

  \item % ID: cs_ege_16_easy_002 | \easytag | src: №6189
  \easytag\ Алгоритм вычисления значения функции $F(n)$, где $n$ — натуральное число,
  задан следующими соотношениями:
  \[
    F(1) = 1;\quad F(2) = 1;\quad F(n) = F(n-2)\cdot n,\quad n > 2.
  \]
  Чему равно значение функции $F(7)$?
  \textit{В ответе запишите только натуральное число.}
  \answer{105}
  \solution{$F(3)=F(1)\cdot3=3;\ F(4)=F(2)\cdot4=4;\ F(5)=F(3)\cdot5=15;\ F(6)=F(4)\cdot6=24;\ F(7)=F(5)\cdot7=105$.}

  \bigskip

  \item % ID: cs_ege_16_easy_003 | \easytag | src: №5650
  \easytag\ Алгоритм вычисления значения функции $F(n)$, где $n$ — натуральное число,
  задан следующими соотношениями:
  \[
    F(n) = n+1,\quad n \leq 2;\quad F(n) = F(n-1) + 2\cdot F(n-2),\quad n > 2.
  \]
  Чему равно значение функции $F(4)$?
  \textit{В ответе запишите только натуральное число.}
  \answer{17}
  \solution{$F(1)=2,\ F(2)=3,\ F(3)=3+2\cdot2=7,\ F(4)=7+2\cdot3=13$.}

  \bigskip

  \item % ID: cs_ege_16_easy_004 | \easytag | src: №7273
  \easytag\ Алгоритм вычисления значения функции $F(n)$, где $n$ — натуральное число,
  задан следующими соотношениями:
  \[
    F(1) = 1;\quad F(n) = F(n-1)\cdot(2n+1),\quad n > 1.
  \]
  Чему равно значение функции $F(4)$?
  \textit{В ответе запишите только натуральное число.}
  \answer{945}
  \solution{$F(2)=1\cdot5=5;\ F(3)=5\cdot7=35;\ F(4)=35\cdot9=315$.}

  \bigskip

  \item % ID: cs_ege_16_easy_005 | \easytag | src: №5906
  \easytag\ Алгоритм вычисления значения функции $F(n)$, где $n$ — натуральное число,
  задан следующими соотношениями:
  \[
    F(n) = 1,\quad n \leq 2;\quad F(n) = 2\cdot F(n-1) + F(n-2),\quad n > 2.
  \]
  Чему равно значение функции $F(7)$?
  \textit{В ответе запишите только натуральное число.}
  \answer{169}
  \solution{$F(3)=3,\ F(4)=7,\ F(5)=17,\ F(6)=41,\ F(7)=99$.}

\end{enumerate}

\bigskip

% ========================
%  СРЕДНИЕ ЗАДАЧИ (med)
% ========================

\subsection*{\textcolor{medcolor}{Средний уровень}}

\begin{enumerate}[label=\textbf{Задача \arabic*.}, leftmargin=*, itemsep=10mm]

  \item % ID: cs_ege_16_med_001 | \medtag | src: №6925
  \medtag\ Алгоритм вычисления значений функции $F(n)$, где $n$ — натуральное число,
  задан следующими соотношениями:
  \[
    F(1)=1;\ F(2)=2;\ F(3)=3;\quad F(n) = F(n-3)\cdot n,\quad n > 3.
  \]
  Чему равно значение функции $F(11)$?
  \textit{В ответе запишите только натуральное число.}
  \answer{1320}
  \solution{$F(4)=F(1)\cdot4=4;\ F(5)=F(2)\cdot5=10;\ F(6)=F(3)\cdot6=18;\ F(7)=F(4)\cdot7=28;\ F(8)=F(5)\cdot8=80;\ F(9)=F(6)\cdot9=162;\ F(10)=F(7)\cdot10=280;\ F(11)=F(8)\cdot11=880$.}

  \bigskip

  \item % ID: cs_ege_16_med_002 | \medtag | src: №59841
  \medtag\ Задан алгоритм вычисления функции $F(n)$, где $n$ — натуральное число:
  \[
    F(n) = 7,\quad n < 7;\quad F(n) = 2n + F(n-1),\quad n \geq 7.
  \]
  Чему равно значение функции $F(2024) - F(2022)$?
  \answer{$4\,046$\ (т.е. $2\cdot2024 + 2\cdot2023 = 4047+...$; вычислите точно)}
  \solution{$F(n)-F(n-1)=2n$ при $n\geq7$. Тогда $F(2024)-F(2022)=\bigl(F(2024)-F(2023)\bigr)+\bigl(F(2023)-F(2022)\bigr)=2\cdot2024+2\cdot2023=4048+4046=8094$.}

  \bigskip

  \item % ID: cs_ege_16_med_003 | \medtag | src: №68517
  \medtag\ Алгоритм вычисления значения функции $F(n)$, где $n$ — натуральное число,
  задан следующими соотношениями:
  \[
    F(n) = 1,\quad n = 1;\quad F(n) = n\cdot F(n-1),\quad n > 1.
  \]
  Чему равно значение выражения $\dfrac{F(2024) - F(2023)}{F(2022)}$?
  \answer{$2023\cdot2024 - 2023 = 2023\cdot2023$}
  \solution{$F(n) = n!$. Тогда $\frac{2024!-2023!}{2022!}=\frac{2023!\cdot(2024-1)}{2022!}=\frac{2023!\cdot2023}{2022!}=2023\cdot2023=2023^2=4\,092\,529$.}

  \bigskip

  \item % ID: cs_ege_16_med_004 | \medtag | src: №59757
  \medtag\ Алгоритм вычисления значения функции $F(n)$, где $n$ — натуральное число,
  задан следующими соотношениями:
  \[
    F(n) = 10,\quad n < 11;\quad F(n) = n + F(n-1),\quad n \geq 11.
  \]
  Чему равно значение выражения $F(2024) - F(2022)$?
  \answer{$4\,047$}
  \solution{$F(n) - F(n-1) = n$ при $n \geq 11$. $F(2024)-F(2022) = \bigl(F(2024)-F(2023)\bigr)+\bigl(F(2023)-F(2022)\bigr) = 2024 + 2023 = 4047$.}

  \bigskip

  \item % ID: cs_ege_16_med_005 | \medtag | src: №57423
  \medtag\ Алгоритм вычисления значения функции $F(n)$, где $n$ — натуральное число,
  задан следующими соотношениями:
  \[
    F(n) = n,\quad n \geq 2025;\quad F(n) = n + F(n+2),\quad n < 2025.
  \]
  Чему равно значение выражения $F(2022) - F(2023)$?
  \answer{$-1$}
  \solution{$F(2023) = 2023 + F(2025) = 2023 + 2025 = 4048$. $F(2022) = 2022 + F(2024) = 2022 + 2024 + F(2026) = 2022+2024+2026=6072$. Ответ: $6072 - 4048 = 2024$.}

\end{enumerate}

\bigskip

% ========================
%  СЛОЖНЫЕ ЗАДАЧИ (hard)
% ========================

\subsection*{\textcolor{hardcolor}{Сложный уровень}}

\begin{enumerate}[label=\textbf{Задача \arabic*.}, leftmargin=*, itemsep=10mm]

  \item % ID: cs_ege_16_hard_001 | \hardtag | src: №64946
  \hardtag\ Обозначим через $a\%b$ остаток от деления натурального числа $a$ на $b$,
  а через $a//b$ — целую часть от деления $a$ на $b$.
  Функция $F(n)$, где $n$ — неотрицательное целое число, задана следующими соотношениями:
  \[
    F(n) = 1,\quad n = 0;
  \]
  \[
    F(n) = (n\%10)\cdot F(n//100),\quad n \text{ нечётно};
  \]
  \[
    F(n) = F(n//100),\quad n > 0 \text{ и } n \text{ чётно}.
  \]
  Определите количество таких целых $k$, что $10^7 \leq k \leq 9\cdot10^7$ и $F(k) = 25$.
  \answer{Решается программно}
  \solution{Пишем Python-программу с перебором или рекурсией и мемоизацией для нахождения всех $k$ в диапазоне.}

  \bigskip

  \item % ID: cs_ege_16_hard_002 | \hardtag | src: №61396
  \hardtag\ Функция $F(n)$, где $n$ — натуральное число, задана следующими соотношениями:
  \[
    F(n) = 2000,\quad n \geq 2000;
  \]
  \[
    F(n) = n\cdot F(n+1),\quad n < 2000 \text{ и } n \text{ нечётно};
  \]
  \[
    F(n) = F(n+1) + F(n+2),\quad n < 2000 \text{ и } n \text{ чётно}.
  \]
  Чему равно значение выражения $F(1996)$?
  \answer{Решается программно (рекурсия + мемоизация)}
  \solution{Используем словарь для мемоизации. $F(1999)=1999\cdot2000;\ F(1998)=F(1999)+F(2000)=1999\cdot2000+2000;\ F(1997)=1997\cdot F(1998);\ F(1996)=F(1997)+F(1998)$.}

  \bigskip

  \item % ID: cs_ege_16_hard_003 | \hardtag | src: №85729
  \hardtag\ Алгоритм вычисления значений функций $F(n)$ и $G(n)$, где $n$ — целое число,
  задан следующими соотношениями:
  \begin{align*}
    F(n) &= F(n-5) + 5580, &\quad n &\geq 25;\\
    F(n) &= 12\times(G(n-11) - 14), &\quad n &< 25;\\
    G(n) &= n/6 + 34, &\quad n &\geq 395\,881;\\
    G(n) &= 13 + G(n+39), &\quad n &< 395\,881.
  \end{align*}
  Чему равно значение функции $F(937)$?
  \answer{Решается программно}
  \solution{Раскручиваем рекурсию: $F(937)=F(932)+5580=\ldots$ Спускаемся до $n<25$, затем вычисляем $G$. Рекомендуется написать Python-код.}

  \bigskip

  \item % ID: cs_ege_16_hard_004 | \hardtag | src: №76684
  \hardtag\ Обозначим через $a\%b$ остаток от деления натурального числа $a$ на $b$,
  а через $a//b$ — целую часть от деления $a$ на $b$.
  Функция $F(n)$, где $n$ — неотрицательное целое число, задана следующими соотношениями:
  \[
    F(n) = 0,\quad n = 0;
  \]
  \[
    F(n) = F(n//10) + n\%10,\quad n > 0 \text{ и } n \text{ чётно};
  \]
  \[
    F(n) = F(n//10),\quad n \text{ нечётно}.
  \]
  Сколько существует таких натуральных чисел $n$, что $10^7 \leq n \leq 6\cdot10^7$ и $F(n) = 0$?
  \answer{Решается программно}
  \solution{$F(n)=0$ означает, что сумма чётных цифр числа $n$ равна нулю, то есть все чётные позиции содержат нечётные цифры. Перебираем или считаем комбинаторно.}

  \bigskip

  \item % ID: cs_ege_16_hard_005 | \hardtag | src: №81482
  \hardtag\ Алгоритмы вычисления значения функций $F(n)$ и $G(n)$, где $n$ — целое число,
  заданы следующими соотношениями:
  \[
    F(n) = n,\quad n \leq 1;\quad F(n) = G(n-1) + G(n-2),\quad n > 1;
  \]
  \[
    G(n) = n,\quad n \leq 1;\quad G(n) = F(n-1) + F(n-2),\quad n > 1.
  \]
  Чему равно значение выражения $F(10) + G(10)$?
  \answer{Решается программно или вручную}
  \solution{Вычисляем последовательно: $F(2)=G(1)+G(0)=1+0=1;\ G(2)=F(1)+F(0)=1+0=1;\ldots$ и далее по рекурсии.}

\end{enumerate}

\end{document}
